
\section*{Lista de Exercício – Seleção de Materiais - Prova P1}

\begin{itemize}
    \item \textbf{Professor}: Dr. Márcio Wagner
    \item \textbf{Disciplina}: Seleção de Materiais
    \item \textbf{Data de Entrega}: 12/02/2025
    \item \textbf{Aluno}: Alan Henrique Pereira Miranda
    \item \textbf{Matrícula}: 202102140072
\end{itemize}

\section{Questão 1}
\subsection*{a) Fluxograma de Projeto}
O fluxograma de projeto é uma representação visual das etapas envolvidas no desenvolvimento de um produto ou sistema. Ele inclui fases como identificação do problema, definição de requisitos, concepção, seleção de materiais, modelagem, prototipagem e testes.

\subsection*{b) Sistema Técnico e suas Subdivisões}
Um sistema técnico é um conjunto de elementos que trabalham juntos para realizar uma função específica. Suas subdivisões incluem:
\begin{itemize}
    \item \textbf{Entrada}: matérias-primas, informação, energia.
    \item \textbf{Processamento}: transformação de insumos.
    \item \textbf{Saída}: produto final, serviço ou energia gerada.
    \item \textbf{Controle}: ajustes e feedbacks para manutenção do sistema.
\end{itemize}

\subsection*{c) Estrutura da Função}
Define a relação entre os elementos de um sistema e como esses elementos interagem para cumprir a função desejada. É composta por funções principais e secundárias.

\subsection*{d) Caminho Convoluto do Projeto}
Refere-se às iterativas e complexas interações entre diferentes etapas do projeto, que podem exigir revisões contínuas.

\subsection{*e) Problema Central de Seleção de Materiais}
A seleção envolve a escolha do material ideal considerando função, processo de fabricação, forma e custo.

\section{Famílias de Materiais e seus Principais Membros}
\begin{itemize}
    \item \textbf{Metais}: Aço, alumínio, cobre, titânio.
    \item \textbf{Cerâmicas}: Alumina, zircônia, carbeto de silício.
    \item \textbf{Polímeros}: Polietileno, polipropileno, nylon.
    \item \textbf{Compósitos}: Fibra de carbono, fibra de vidro.
    \item \textbf{Elastômeros}: Borracha natural, neoprene, silicone.
\end{itemize}

\section{Propriedades dos Materiais}
\begin{itemize}
    \item \textbf{Cerâmicas}: Alta dureza, baixa tenacidade, resistência ao calor.
    \item \textbf{Vidros}: Transparência, fragilidade, resistência a ataques químicos.
    \item \textbf{Polímeros}: Flexibilidade, baixa densidade, boa moldabilidade.
    \item \textbf{Elastômeros}: Alta elasticidade, resistência ao impacto.
    \item \textbf{Metais}: Alta resistência mecânica, ductilidade, condutividade térmica.
    \item \textbf{Híbridos}: Combina propriedades de diferentes materiais, como rigidez e leveza.
\end{itemize}

\section{Tipos de Informações Necessárias aos Projetos}
\begin{itemize}
    \item Propriedades mecânicas e físicas.
    \item Processabilidade.
    \item Durabilidade e custos.
\end{itemize}

\section{Materiais Quanto à Ductilidade}
\subsection*{a) Dúcteis}
Sofrem grandes deformações plásticas antes da fratura.

\subsection*{b) Polímeros Dúcteis}
Apresentam deformabilidade acentuada, podendo esticar significativamente antes de romper.

\subsection*{c) Cerâmicas Frágeis}
Rompem com pouca deformação plástica.

\section{O que é Limite de Fadiga?}
A tensão abaixo da qual um material pode suportar um número infinito de ciclos de carga sem falhar.

\section{O que é Dureza?}
Resistência à penetração ou ao desgaste. Principais escalas: Brinell, Rockwell, Vickers.

\section{O que é Tenacidade à Fratura?}
Capacidade de um material resistir à propagação de trincas.

\section{Condutividade Térmica}
Medida da capacidade de transferência de calor. Governada pela equação de Fourier.

\section{Expansão Térmica}
Variação dimensional de um material com a temperatura, descrita pelas equações:

\begin{itemize}
    \item \textbf{Linear}: \( \Delta L = \alpha L_0 \Delta T \)
    \item \textbf{Superficial}: \( \Delta A = 2\alpha A_0 \Delta T \)
    \item \textbf{Volumétrica}: \( \Delta V = 3\alpha V_0 \Delta T \)
\end{itemize}

\section{Como são obtidos os valores de resistência dos materiais?}
Por meio de ensaios mecânicos como tração, compressão e impacto.

\section{Diagrama de Barras Únicas}
Representação gráfica de uma propriedade de vários materiais.

\section{Diagramas de Propriedades Múltiplas}
Gráfico com duas propriedades distintas para seleção de materiais.

\section{Explicação dos Diagramas}
\begin{itemize}
    \item \( E-\rho \) (Módulo de Young vs. Densidade)
    \item \( R-\rho \) (Resistência vs. Densidade)
    \item \( E^*/R^* \) (Módulo Específico vs. Resistência Específica)
    \item Tenacidade vs. Módulo
    \item Tenacidade à Fratura vs. Resistência
\end{itemize}

\section{Por que Conectar Material e Função?}
Para garantir que o material atenda às demandas do projeto com eficiência.

\section{Estratégia para Seleção de Materiais}
Definir requisitos, classificar opções e otimizar a escolha.

\section{Passos: Tradução, Triagem, Classificação e Documentação}
Transformar requisitos do projeto em critérios técnicos e realizar a seleção.

\section{O que é Função, Restrição, Objetivo e Variável Livre?}
\begin{itemize}
    \item \textbf{Função}: Propósito do componente.
    \item \textbf{Restrição}: Limites obrigatórios.
    \item \textbf{Objetivo}: O que otimizar.
    \item \textbf{Variável Livre}: Parâmetros ajustáveis.
\end{itemize}
